% !PRINTED 47
\weberpage{47}
%<<CONT

\[
M_1=\Delta_0=f(1)-f(0).
\]
Put
\begin{equation}
\begin{aligned}
 \Delta_{x+1}-\Delta_x&=\Delta_x',\\
 \Delta_{x+1}'-\Delta_x'&=\Delta_x'',\\
 &\cdots\\
 \Delta_{x+1}^{(n-2)}-\Delta_x^{(n-2)}&=\Delta_x^{(n-1)}.
\end{aligned}
\tag{6}
\end{equation}
then, accordingly,
\[
 M_0=f(0),\quad M_1=\Delta_0,
 \quad M_2=\Delta_0',\quad\ldots,\quad M_n=\Delta_0^{(n-1)},
\]
and formula (3) gives
\begin{equation}
 f(x)=f(0)+\Delta_0B_1^{(x)}+\Delta_0'B_2^{(x)}+\Delta_0''B_3^{(x)}+\cdots+\Delta_0^{(n-1)}B_n^{(x)}.
\tag{7}
\end{equation}
Likewise the functions $\Delta_x,\Delta_x',\Delta_x'',\ldots$ can be expressed:
\begin{equation}
\begin{aligned}
 \Delta_x&=\Delta_0+\Delta_0'B_1^{(x)}+\Delta_0''B_2^{(x)}+\cdots+\Delta_0^{(n-1)}B_{n-1}^{(x)},\\
 \Delta_x'&=\Delta_0'+\Delta_0''B_1^{(x)}+\cdots+\Delta_0^{(n-1)}B_{n-2}^{(x)}.
\end{aligned}
\tag{8}
\end{equation}
and so on.

The $\Delta_x,\Delta_x',\Delta_x'',\ldots,\Delta_x^{(n-1)}$ are whole rational functions of degrees $n-1,n-2,\ldots,0$; the last of them is therefore constant. In order to represent them all, one needs only the values
\[
 f(0),\quad \Delta_0,\quad \Delta_0',\quad \Delta_0'',\quad\ldots,\quad \Delta_0^{(n-1)},
\]
which one computes most easily by arranging a table. For the case $n=3$, for example, the table has the form:
\[
\begin{array}{c|c|c|c}
f(0) &&&\\
f(1)&\Delta_0&&\\
f(2)&\Delta_1&\Delta_0'&\\
f(3)&\Delta_2&\Delta_1'&\Delta_0''
\end{array}
\]
When $f(0),f(1),\ldots$ are given, the table is filled by simple subtractions.

\begin{center}
{\large \S.\ 11.}

\vspace{0.5em}
{\large\gesperrt{Arithmetic Series of Higher Order.}}
\end{center}

The foregoing is the foundation for the theory of arithmetic series of higher order.

We denote by
\begin{equation}
 u_0,u_1,u_2,u_3,\ldots
\tag{1}
\end{equation}
any infinite series of quantities and by
%CONT>>
